\section{Attack Implementations}
\label{sec:attacks}

\subsection{AgentPoison}

\agentpoison{} \citep{chen2024agentpoison} optimizes a trigger sequence
$\trigger^* = (t_1, \ldots, t_\ell)$ such that for any victim query
$q \in \mathcal{Q}_v$, the triggered embedding $E(q \oplus \trigger)$ is
close to the adversarial passage embedding $E(\passage)$:
\begin{equation}
  \trigger^* = \arg\max_\trigger\; \EE_{q \sim \mathcal{Q}_v}
  \left[\cos\!\left(E(q \oplus \trigger),\, E(\passage)\right)\right].
  \label{eq:agentpoison_obj}
\end{equation}

\paragraph{Centroid passage (our implementation).}
The full HotFlip-style~\citep{ebrahimi2018hotflip} gradient optimization
over DPR embeddings requires GPU-resident gradients.  We approximate it with
vocabulary coordinate descent over the all-MiniLM-L6-v2 embedding space
(Phase 10) and a \emph{centroid-targeting adversarial passage}:
\begin{equation}
  \passage = \mathrm{CentroidPassage}(\mathcal{Q}_v)
  = \text{passage covering vocabulary}\; \bigcup_{q \in \mathcal{Q}_v} \mathrm{KeyTerms}(q).
  \label{eq:centroid}
\end{equation}
$\mathrm{KeyTerms}$ extracts content words (stopword-filtered, frequency-ranked).
This places $\passage$ near the centroid of victim query embeddings, achieving
moderate-to-high cosine similarity with \emph{every} victim query rather than
maximum similarity with a single query.

\paragraph{Triggered-query evaluation.}
A critical evaluation detail: the attacker controls query content at test
time (e.g., via a shared trigger deployed to compromised clients).  The
correct evaluation protocol (Chen et al., Algorithm~2) issues triggered
queries $q \oplus \trigger$, not plain queries $q$.  We correct this in our
framework and discuss the resulting performance improvement in
Section~\ref{sec:experiments}.

\paragraph{Modelled $\asra$.}
Conditional action execution probability is modelled from paper-reported
values: $\asra \sim \mathcal{N}(0.68, 0.06^2)$ per retrieval event.

\subsection{MINJA}

\minja{} \citep{dong2025minja} uses a \emph{progressive-shortening indication
protocol}: the attacker sends $T$ crafted queries per target entry, each
progressively shorter to remain below detection thresholds.  The
injection success rate is:
\begin{equation}
  \isr = \frac{|\{\text{entries successfully stored}\}|}{n_\text{poison}},
  \qquad
  \Pr[\text{turn } t \text{ succeeds}] = p_0 \cdot \max(0.5,\, 1 - \lambda t),
\end{equation}
where $p_0 = 0.98$ (full indication probability) and $\lambda = 0.10$
(shortening rate).  With $T=3$ turns, ISR $\approx 0.982$, reproducing the
paper result.

Adversarial passages use \emph{bridging steps} that connect the victim query
topic to the malicious goal via a procedurally correct reasoning chain.
Modelled $\asra \sim \mathcal{N}(0.76, 0.05^2)$ (highest of the three
attacks due to convincing bridging framing).

\subsection{InjecMEM}

\injecmem{} \citep{injecmem2026} employs a retriever-agnostic anchor: a
passage covering broad agent task vocabulary (meetings, tasks, calendar,
security, infrastructure, etc.) such that many query types partially match it.
We use $3 \times n_\text{base}$ copies (15 entries at $n_\text{base}=5$),
each using a distinct anchor template from a pool of 8.  Modelled
$\asra \sim \mathcal{N}(0.57, 0.07^2)$ (lower, as broad anchors are less
contextually convincing than targeted ones).

\subsection{Evaluation Setup}

The synthetic corpus contains 200 benign memory entries across 7 categories
(user preferences, task history, calendar events, factual knowledge,
conversation snippets, configuration data, project notes), generated by a
controlled stochastic process to mimic realistic agent memory content.
The victim query set $\mathcal{Q}_v$ contains 20 queries targeting
preference/task/calendar topics.  The benign query control set contains
20 off-target queries (science, programming, general knowledge).
